\chapter{命题逻辑：语法与语义}
\begin{qlnote}{阅读参照}Enderton, \emph{Logic}, Chapter One.\end{qlnote}

\section{递归定义公式}
\begin{definition}{Sentential formulas}
给定命题字母集合，取 \(\neg,\to\) 为基本联结词。每个命题字母是公式；
若 \(\varphi,\psi\) 是公式，则 \((\neg\varphi)\) 和 \((\varphi\to\psi)\)
也是公式；除此之外没有公式。括号是语法的一部分。
\end{definition}
这一定义同时给出\textbf{结构归纳}：先证明命题字母满足性质，再证明两个构造步骤保留性质，
便可得到所有公式满足该性质。公式是一棵有限语法树；它的语义通过树递归计算。

\section{真值赋值}
\begin{definition}{Truth assignment}
赋值 \(v\) 把每个命题字母映到 \(\{0,1\}\)，并唯一递归扩张到公式：
\[
 v(\neg\varphi)=1-v(\varphi),\qquad
 v(\varphi\to\psi)=0\ \Longleftrightarrow\
 v(\varphi)=1\text{ 且 }v(\psi)=0.
\]
\end{definition}
\(\varphi\land\psi\)、\(\varphi\lor\psi\) 和 \(\varphi\leftrightarrow\psi\)
可作为缩写引入，也可作为带真值规则的基本联结词。前后必须保持同一种约定。

\begin{definition}{Satisfiability and semantic consequence}
若存在赋值使一组公式 \(\Gamma\) 全为真，则 \(\Gamma\) \textbf{可满足}。
\(\Gamma\models\varphi\) 表示每个使 \(\Gamma\) 全为真的赋值也使 \(\varphi\) 为真。
若 \(\varnothing\models\varphi\)，则 \(\varphi\) 是\textbf{重言式}。
若 \(\models\varphi\leftrightarrow\psi\)，则两公式\textbf{逻辑等价}。
\end{definition}

\begin{theorem}{Modus ponens preserves truth}
\(\{\varphi,\varphi\to\psi\}\models\psi\)。
\end{theorem}
\begin{proof}
取使两个前提为真的赋值。若 \(v(\psi)=0\)，由于 \(v(\varphi)=1\)，
条件公式的真值规则给出 \(v(\varphi\to\psi)=0\)，矛盾。
\end{proof}

\section{充分、必要与逆否}
\(P\to Q\) 中，\(P\) 是 \(Q\) 的充分条件，\(Q\) 是 \(P\) 的必要条件。
它与 \(\neg Q\to\neg P\) 逻辑等价。逆命题 \(Q\to P\) 和否命题
\(\neg P\to\neg Q\) 互相等价，却一般不由原命题推出。
赋值 \(v(P)=0,v(Q)=1\) 就使原命题真、逆命题假。

\begin{example}[真值表能检查什么]
一个只含有限多个命题字母的公式，可由有限真值表判定是否为重言式。
例如 \(P\lor\neg P\) 在经典语义下总为真。
这项算法检查的是命题联结结构；它不会检查「连续函数一定可导」这样的数学前提是否成立。
\end{example}

\section{练习}
证明 De Morgan 等价式；为肯定后件 \((P\to Q),Q\therefore P\) 找反例；
用结构归纳证明每个公式只含有限多个命题字母。阅读原书时补足唯一可读性与命题紧致性的证明。
